Nesting tikzpictures
    [ i.e. create another tikzpicture inside the content/text of a node (\node(){\tikz{...}}) ]
is actually discouraged and thus not guaranteed to work fully as expected.
E.g. there may occur shifts in objects, coordinates are not properly accessible, properties inherit, scaling and other transformations may mess up.
For simple stuff, it may satisfy you, for more complex things, I suggest one of the following.

    | Sidenote:
        [ At least for such stuff, draw/create/compile/render/include your tikzpictures as standalone-docs. That way, the pics are nicely isolated from the actual document. That saves resources, as for instance variables and registers (refer to savebox) are only instantiated inside the tikzpic and vanish afterwards. Also you can utilize functions like "buildnew", so your compilation time greatly improves since pics do not have to be compiled repeatedly as long as they don't change. ]
        [ I have to say: I do pretty much ALWAYS use standalone for tikzpics and then include them. It's pretty neat. ]




#####################
1) Use saveboxes to then safely nest tikzpictures
    (-> Coordinates will be remembered. The way [remember picture] works is to save the origin location in the aux file, so you have to run the code twice.)
    (It should be noted that by placing the picture inside a node you are centering it (default) at the node location (origin).)
!(!
    When using this together with standalone, you have to take one precaution, that is to "disable" the \newpage-cmd locally. This is due to a \newpage that "standalone" apparently attaches to tikzpicture. As you cannot have a new page in an sbox, it suffices to just disable this cmd locally with "\let\newpage\relax". See below for an example on that
!)!
(-> Be careful to declare the \sbox outside the actual tikzpicture or you run into other problems, like not working LaTeX directives and formatting or surplus whitespaces)
#####################
\documentclass{article}
\thispagestyle{empty}
\usepackage{graphicx}
\usepackage{tikz}
\begin{document}
\begin{figure}
\sbox0{\begin{tikzpicture}[remember picture]
    \coordinate(v)at(3, 4);
    \node(t)at(v){v};
    \draw[->](1,2)--(v);
    \end{tikzpicture}
}
\begin{tikzpicture}[remember picture]
    \node(inset){\usebox0};
    \node(t2)[blue]at(v){v};
    \draw[->,red](5,-2)--(v);
\end{tikzpicture}
\end{figure}
\end{document}
% = = = = = = = = = = =
% -- Disabling the newpage-cmd
\newsavebox{\someSBox}%
\sbox{\someSBox}{\let\newpage\relax\begin{tikzpicture}[remember picture]%
    \begin{scope}[dashed,blue]
        \node(){test};
    \end{scope}
\end{tikzpicture}}%
\usebox{\someSBox}%




#####################
2) Instead of nesting tikzpictures, use a "scope" with "local bounding box"
    (Gives you "wrapper" around the scope for e.g. attaching stuff to anchors, like drawing a connection onto the "collection of objects")
#####################
\documentclass{article}
\thispagestyle{empty}
\usepackage{tikz}
\usetikzlibrary{fit}
\begin{document}
\begin{figure}
\begin{tikzpicture}
\begin{scope}[local bounding box=inset]
    \coordinate(v)at(3, 4);
    \node(t)at(v){v};
    \draw[->](1,2)--(v);
\end{scope}
\node(t2)[blue]at(v){v};
\draw[->,red](5,-2)--(v);

\draw[red](inset.east)--(inset.west);
\end{tikzpicture}
\end{figure}
\end{document}




#####################
3) In certain cases, for certain features (accessing coordinates of nested picture) it might work to append the "remember picture" property
    (I consider this rather clumsy)
#####################
\documentclass{article}
\thispagestyle{empty}
\usepackage{tikz}
\begin{document}
\tikzstyle{every picture}+=[remember picture]
\begin{tikzpicture}
    \node(inset){
        \begin{tikzpicture}[remember picture]
            \coordinate(v)at(3,4);
            \node(t)at(v){v};
            \draw[->](1,2)--(v);
        \end{tikzpicture}
    };
    \node(t2)[blue]at(v){v};
    \draw[->,red](5,-2)--(v);
\end{tikzpicture}
\end{document}